En este Trabajo Fin de Grado se desarrolla el cálculo y diseño de una línea aérea de alta tensión de \SI{20}{\kilo\volt}, de simple circuito y disposición simplex, destinada a la evacuación de la energía generada por una planta fotovoltaica de \SI{6.5}{\mega\watt}. El proyecto se ha estructurado siguiendo los apartados propios de un Proyecto Tipo de Líneas Aéreas de Alta Tensión con tensión nominal inferior a \SI{36}{\kilo\volt}, de modo que tanto el índice como el contenido y la documentación técnica se ajustan a los criterios habituales de este tipo de instalaciones.

El trazado de la línea se sitúa en la provincia de Alicante, entre los municipios de Benidorm y L’Albir, zona caracterizada por una orografía y condiciones ambientales que condicionan el diseño mecánico, eléctrico y de seguridad de la instalación. El estudio incluye la definición del trazado en planta y perfil, la selección de conductores de aluminio-acero, apoyos, armados, cadenas de aisladores y sistemas de protección, así como el cumplimiento de las distancias reglamentarias y prescripciones ambientales y de seguridad.

El diseño de la línea se ha realizado mediante el software IMEDEXSA. Paralelamente, todos los cálculos eléctricos, mecánicos, de cimentaciones y de puesta a tierra han sido programados y desarrollados de forma independiente en MATHCAD siguiendo los procedimientos y criterios establecidos por la normativa vigente., con el objetivo de verificar la coherencia y validez de los resultados obtenidos con el software de diseño. Este enfoque permite contrastar los resultados y reforzar la comprensión del proceso de cálculo de una línea aérea de alta tensión.

Se desarrollan los cálculos eléctricos, incluyendo intensidades máximas admisibles, caída de tensión, pérdidas de potencia y rendimiento de la línea, así como los cálculos mecánicos de tracciones y flechas bajo distintas hipótesis de carga. Asimismo, se estudian las cimentaciones monobloque de los apoyos, aplicando el método de cálculo correspondiente, y se diseña el sistema de puesta a tierra, verificando el cumplimiento de las tensiones de paso y contacto admisibles para garantizar la seguridad de las personas.

Finalmente, el proyecto se completa con los planos constructivos, el pliego de condiciones técnicas, el presupuesto y el estudio de seguridad y salud, obteniendo como resultado una solución técnica viable, coherente con la normativa y representativa de un proyecto real de línea aérea de media tensión para evacuación de energía renovable.
\clave{Línea Aérea de \SI{20}{\kilo\volt}, Simple Circuito, Simplex, Alicante, IMEDEXSA, MATHCAD, Planos, Pliego de Condiciones, Presupuesto, Estudio de Seguridad y Salud, Cimentaciones, Puesta a Tierra, Energía Renovable}